\reading{CR-2}{Governance}{DEFER}{1}

\begin{crkeybox}[Learning objectives for this reading]
\begin{enumerate}[label=\textbf{\alph*.},leftmargin=2em]
\item Define risk management responsibilities in an organization and explain the three lines of defense framework for effective risk management and control.
\item Explain the processes that lead to risk taking including credit origination, credit risk assessment, and credit approval processes.
\item Discuss the following key principles underlying best practice for the governance system of credit risk: Guidelines, Skills, Limits, and Oversight.
\item Describe the most common parameters of a credit-sensitive transaction.
\item Describe the roles of the credit committee in an organization.
\end{enumerate}
\end{crkeybox}

% =====================================================================
\lo{CR-2 a}{Define risk management responsibilities in an organization and explain the three lines of defense framework for effective risk management and control.}

\begin{center}
\footnotesize
\begin{tabular}{L{2.5cm} L{5.0cm} L{7.1cm}}
\toprule
\thead{Line} & \thead{Who} & \thead{What they do} \\
\midrule
\term{First line} & Business owners &
  Primarily \emph{own} and manage risk. They manage daily risks and ensure
  adherence to guidelines. \\
\term{Second line} & Enterprise risk management, compliance, and legal &
  Monitor and oversee risk. They establish policies and procedures and serve as
  management oversight for the first line. \\
\term{Third line} & Internal audit, external auditors, and special audit
  committees &
  Provide \emph{independent assurance} of the risk management and risk
  monitoring provided by the first and second lines of defense. \\
\bottomrule
\end{tabular}
\end{center}
\figcap{The three lines of defense. The distinction that carries the weight is
own versus oversee versus independently assure, and which function sits in each
line.}

An organization's \term{chief executive officer (CEO)} is responsible for
ensuring that adequate governance guidelines have been established, and that
proper risk governance is in place. The CEO should have the support of his staff
in creating and enforcing robust risk management guidelines.

% =====================================================================
\lo{CR-2 b}{Explain the processes that lead to risk taking including credit origination, credit risk assessment, and credit approval processes.}

Governance ensures organizational functionality and risk management, crucial for
detecting and mitigating risks like individual errors, process breakdowns, and
fraud.

The core processes in risk management include:

\begin{enumerate}
\item \term{Credit origination.} Influences profitability; a focus on volume and
returns may increase risk.
\item \term{Credit risk assessment.} Essential for evaluating the potential risk
of transactions.
\item \term{Credit approval.} Ensures that risks are appropriately mitigated
before proceeding.
\end{enumerate}

Transaction quality determines profitability. Without proper monitoring,
high-risk transactions pursued for increased volume and returns can lead to
substantial losses, necessitating robust risk assessment and approval
mechanisms.

% =====================================================================
\lo{CR-2 c}{Discuss the following key principles underlying best practice for the governance system of credit risk: Guidelines, Skills, Limits, and Oversight.}

Effective governance is based on the following four key principles, which ensure
that goals and objectives are achieved, and risks are mitigated:

\begin{enumerate}
\item \term{Guidelines.} Organizations must institute clear guidelines around
approvals, including approvals for transactions that give rise to credit risk.
\item \term{Skills.} Authority must be delegated to those with proper skills.
\item \term{Limits.} Adequate risk and transaction limits must be set.
\item \term{Oversight.} Employees and functions should be subject to oversight by
qualified and independent people.
\end{enumerate}

\subsection*{1) Guidelines: the framework for action}

Guidelines, such as credit policies or risk standards, serve as the rulebooks for
transaction conduct, emphasizing clarity, brevity, specificity, and
accessibility:

\begin{itemize}
\item \term{Understandable.} Ensure simplicity and clarity in language to
facilitate easy interpretation. Clear, jargon-free language.
\item \term{Concise.} Maintain a reasonable length to enhance user engagement and
comprehension, and to encourage reading.
\item \term{Precise.} Balance simplicity with specificity, incorporating
real-life scenarios for transactional guidance.
\item \term{Accessible.} Ensure easy availability through electronic document
systems, complemented by concise summaries on the company's intranet for swift
access.
\end{itemize}

\textbf{Creation and approval.} Management oversight ensures guidelines contain
appropriate authorizations, typically sponsored by the CFO or CRO and approved by
the Board or its risk committees. Periodic reviews by the Board and risk
committees ensure guidelines remain accurate and up-to-date with current
practices and regulations, with event-triggered reviews for losses.

\textbf{Maintenance and content.} Guidelines are typically managed by the CRO's
office, drafted, approved, and maintained by senior personnel with expertise in
the organization's products and markets. Regular reviews, especially amid
regulatory changes and significant organizational events like mergers, are
essential for guideline effectiveness and compliance.

\textbf{Guideline topics.} Covering the purpose of the guidelines, the
methodology used, the approval flow, handling of new products and markets, the
review and update process, and consequences for non-compliance ensures
comprehensive guideline coverage.

\begin{crtrapbox}[Breach and noncompliance]
Breaches are rare and may result in employment termination, \emph{with some
exceptions for unintentional breaches} such as credit limit breaches caused by
unexpected currency fluctuations. Central databases for transaction exposure and
credit parameters enable pre-trade checks, preventing unauthorized trades and
ensuring compliance with guidelines.
\end{crtrapbox}

\subsection*{2) Skills: competence in execution}

\begin{itemize}
\item \term{Approval.} The CRO and the Board cannot approve all transactions;
authority is delegated down. Risk management acts in an advisory role, focusing
on guidance rather than direct approval.
\item \term{Authority levels.} Based on risk: simpler transactions require
lower-level approval, while complex or high-risk ones need senior management.
\end{itemize}

\subsection*{3) Limits: the boundaries of risk}

Limits, also known as credit lines, serve as benchmarks for the maximum
acceptable loss an organization is willing to bear, often specified in absolute
dollar terms. This includes both individual and aggregate limits.

For instance, a company might set an overarching absolute risk limit, such as
\$75 million, while also assigning smaller limits to specific counterparties,
industries, or geographic regions.

These limits are established through a blend of risk assessments, management's
experience-driven intuition, and external factors like regulatory requirements
and shareholder preferences. Moreover, individual business units may prefer
preapproved limits to expedite client interactions and streamline decision-making
processes.

\subsection*{4) Oversight: vigilance and independence}

\begin{enumerate}
\item \term{Independence.} Risk management should be separate from profit
centers and have compensation independent of business performance. Reporting
structure: risk management reports to the CRO, who reports to the CEO and has
access to the risk and audit committees.
\item \term{Strong qualifications.} Risk managers need a deep understanding of
the business, profit motivations, and transaction structures.
\item \term{Closeness to the business.} While independent, risk managers should
collaborate and attend client meetings when necessary. They should maintain a
clear final position based on risk limits and avoid emotional or client
pressure, balance business needs with protecting the organization's interests,
and facilitate deal closures by helping originators succeed.
\item \term{Open-mindedness.} Be receptive to new information and willing to
adapt when necessary, facilitating deal closures by balancing business needs
with risk management.
\end{enumerate}

\begin{crkeybox}[The risk manager's role]
Balance business needs with protecting the organization. Risk managers should
\emph{facilitate} deal closures, not block them.
\end{crkeybox}

% =====================================================================
\lo{CR-2 d}{Describe the most common parameters of a credit-sensitive transaction.}

Credit-sensitive transactions can be described by three parameters. Credit risk
boils down to these three key factors:

\begin{enumerate}[label=\Alph*.]
\item \term{Amount of exposure} --- money on the line. Transactions can expose
organizations to losses. Measuring the amount of potential losses is critical.
How much can you lose on a transaction?
\item \term{Credit quality} --- the borrower's reliability. Organizations should
assess the risk of losses from transactions with counterparties. Assessing the
creditworthiness of counterparties is therefore important. How likely is the
other party to repay?
\item \term{Length of exposure (tenor)} --- time to repay. Organizations should
have a good understanding of the period during which they are exposed to
potential losses, which is the length of time until an obligation is due by a
counterparty. How long are you waiting for your money?
\end{enumerate}

% =====================================================================
\lo{CR-2 e}{Describe the roles of the credit committee in an organization.}

The credit committee, comprised of senior executives and senior experts from
across the organization (business, risk, legal, and so on), oversees important or
high-risk transactions with a clear approval process outlined in its charter.

\begin{itemize}
\item It approves high-risk loans based on clear guidelines.
\item It reviews packages prepared by originators in advance, boasting a track
record of decision-making.
\item Committee members, representing key organizational functions, possess deep
expertise and decision-making authority, led by a chair who facilitates
discussions and ensures objectivity.
\item There is a voting process, and meeting minutes are recorded and distributed
promptly after each session.
\end{itemize}
